\documentclass[11pt]{article}

\usepackage[margin=1in]{geometry}
\usepackage{amsmath,amssymb}
\usepackage{booktabs}
\usepackage{tabularx}
\usepackage{array}
\usepackage{enumitem}
\usepackage[T1]{fontenc}
\usepackage{hyperref}

\newcommand{\code}[1]{\texttt{#1}}
\newcommand{\keff}{k_{\mathrm{eff}}}
\newcommand{\Uh}{\hat U}
\newcommand{\MF}{\mathrm{MF}}
\newcommand{\vir}{\mathrm{vir}}
\emergencystretch=3em

\title{Mean-field Dean--Kawasaki versus macroscopic fluctuation theory:\\
virial coefficients, free energy, structure factor and diffusivity}
\author{FHDeX}
\date{}

\begin{document}
\maketitle

\noindent This note explains why the coarse-grained Dean--Kawasaki equation
(DKE) solved in \code{src\_deankow/interaction} and macroscopic fluctuation
theory (MFT) with the exact transport coefficients, solved in
\code{src\_deankow/MFT}, predict different long-wavelength statics and
dynamics for interacting Brownian particles. It also explains why the
nonlinearity and the multiplicative noise of the DKE, and refining its grid,
do not remove the difference. The example throughout is the 1D Gaussian-core
potential of Grabsch, Venturelli and B\'enichou (arXiv:2504.08560,
supplementary material, Sec.~II), whose notation (S$n$) we follow for the
equation numbers quoted from it.

\section{Microscopic model}

$N$ Brownian particles interact through a pair potential $V$:
\begin{equation}\label{eq:langevin}
  \dot x_i = -\mu_0 \sum_{j\neq i} V'(x_i-x_j) + \sqrt{2D_0}\,\xi_i(t),
  \qquad D_0=\mu_0 k_BT,
\end{equation}
with independent unit white noises $\xi_i$. We set $\mu_0=k_BT=D_0=1$ and use
\begin{equation}
  V(r) = \varepsilon\, e^{-r^2/\ell^2}.
\end{equation}
In 1D, the dimensionless control parameters are the coupling $\beta\varepsilon$
and the density $\rho\ell$ (particles per interaction range). The potential has
no hard core, so particles can pass each other.

\section{MFT: exact transport coefficients}

On large scales the density obeys
\begin{equation}\label{eq:mft}
  \partial_t\rho = \partial_x\!\left[D(\rho)\,\partial_x\rho
     + \sqrt{\sigma(\rho)/\Lambda}\;\eta\right],
\end{equation}
where the transport coefficients are fixed by the microscopic model.

\paragraph{Mobility.} $\sigma$ is the linear response of the mean current to a
uniform force $F_0$ on all particles, $\langle j\rangle = \sigma F_0/(2k_BT)$.
A uniform force is removed by a Galilean change of frame, so every particle
drifts at $\mu_0F_0$ whatever the interactions, and
\begin{equation}\label{eq:sigma}
  \sigma(\rho) = 2\mu_0 k_BT\,\rho = 2\rho \qquad\text{(S56)}.
\end{equation}
This is exact and contains no information about $V$.

\paragraph{Diffusivity.} The fluctuation--dissipation relation of MFT,
$2k_BT\,D/\sigma = f''(\rho)$ (S51), with $f$ the equilibrium free energy
density, then gives
\begin{equation}\label{eq:D}
  D(\rho) = \mu_0\,\rho f''(\rho) = \mu_0\,\partial_\rho P(\rho)
  \qquad\text{(S60)},
\end{equation}
using $\partial_\rho P = \rho f''$ for the pressure $P=\rho f'-f$. $D$ is
therefore entirely thermodynamic: it is the inverse compressibility.

\paragraph{Virial expansion.} With
$P/k_BT = \rho + B_2\rho^2 + B_3\rho^3 + B_4\rho^4 + \dots$,
\begin{equation}\label{eq:Dvir}
  D_\vir(\rho) = 1 + 2B_2\rho + 3B_3\rho^2 + 4B_4\rho^3 + \dots
\end{equation}
In 1D, with the Mayer function $f_M(x) = e^{-\beta V(x)}-1$,
\begin{align}
  B_2 &= -\tfrac12\int f_M(x)\,dx, \qquad
  B_3 = -\tfrac13\iint f_M(x)\,f_M(y)\,f_M(y-x)\,dx\,dy, \\
  B_4 &= -\tfrac18\int\!\!\int\!\!\int \big[\,3\,f_{12}f_{23}f_{34}f_{41}
        + 6\,f_{12}f_{23}f_{34}f_{41}f_{13}
        + f_{12}f_{13}f_{14}f_{23}f_{24}f_{34}\big]\,dx_2\,dx_3\,dx_4 ,
\end{align}
with $x_1=0$ and $f_{ij}=f_M(x_i-x_j)$. Table~\ref{tab:B} lists the
coefficients for the Gaussian core, computed by quadrature. (The quadrature
reproduces the Tonks values $B_2=B_3=B_4=1$ for hard rods of unit length to
$1\%$, and is converged in the grid spacing for the Gaussian.) For
$\beta\varepsilon=2$ it reproduces S102,
$D = 1 + 1.97269\rho + 2.12375\rho^2$.

\begin{table}[h]
\centering
\begin{tabular}{ccccc}
\toprule
$\beta\varepsilon$ & $B_2/\ell$ & $B_3/\ell^2$ & $B_4/\ell^3$ & $B_2^\MF/\ell = \sqrt\pi\,\beta\varepsilon/2$ \\
\midrule
1 & 0.6426 & 0.2104 & $-0.0380$ & 0.886 \\
2 & 0.9863 & 0.7079 & 0.2362 & 1.772 \\
3 & 1.1887 & 1.1704 & 0.8390 & 2.659 \\
5 & 1.4124 & 1.8081 & 2.0354 & 4.431 \\
\bottomrule
\end{tabular}
\caption{Virial coefficients of the 1D Gaussian core, and the mean-field value
of $B_2$ (Sec.~\ref{sec:dk}).}
\label{tab:B}
\end{table}

\section{Structure factor and intermediate scattering function}\label{sec:S}

Linearize \eqref{eq:mft} about a uniform $\bar\rho$. A Fourier mode is then an
Ornstein--Uhlenbeck process,
\begin{equation}
  \partial_t\delta\hat\rho_k = -D k^2\,\delta\hat\rho_k
     + ik\sqrt{\sigma/\Lambda}\;\hat\eta_k ,
\end{equation}
which relaxes at rate $Dk^2$ and is driven with strength $\sigma k^2/\Lambda$.
Its stationary variance is $\sigma/(2D\Lambda)$ per unit volume, so with the
normalization $S = \Lambda\langle|\delta\hat\rho_k|^2\rangle/(\bar\rho V)$
(that of $|\hat n_k|^2/N$ for particle counts $n$) the factors of $k^2$ cancel:
\begin{equation}\label{eq:S}
  S(k\to0) = \frac{\sigma(\bar\rho)}{2\bar\rho\,D(\bar\rho)} = \frac{1}{D(\bar\rho)}
  = \frac{k_BT}{\partial_\rho P},
\end{equation}
the compressibility sum rule. The noise injects fluctuations at a rate set by
$\sigma$ and the current removes them at a rate set by $D$; since $\sigma$ is
fixed by \eqref{eq:sigma}, $S$ depends on $D$ alone. The intermediate
scattering function is
\begin{equation}\label{eq:F}
  F(k,\tau) = S(k)\,e^{-D k^2 \tau},
\end{equation}
so the same $D$ sets both the amplitude $F(k,0)=S=1/D$ and the decay rate
$Dk^2$. A wrong $D$ gets both wrong, in a correlated way: a model that
overestimates $D$ predicts too small an $S$ and too fast a decay.

\section{Mean-field Dean--Kawasaki equation}\label{sec:dk}

The coarse-grained DKE solved in \code{src\_deankow/interaction} is
\begin{equation}\label{eq:dk}
  \partial_t\rho = \partial_x\!\left[D_0\partial_x\rho
     + \mu_0\rho\,\partial_x(U*\rho) + \sqrt{2D_0\rho/N}\;\eta\right],
\end{equation}
with $U*\rho$ the convolution with the pair potential. It is nonlinear (the
drift is quadratic in $\rho$) and has multiplicative noise.

\paragraph{Free energy.} \eqref{eq:dk} is a gradient flow with mobility
$\mu_0\rho$ of the functional
\begin{equation}\label{eq:FMF}
  \mathcal F_\MF[\rho] = k_BT\!\int(\rho\ln\rho-\rho)\,dx
     + \frac12\iint \rho(x)\,U(x-y)\,\rho(y)\,dx\,dy ,
\end{equation}
since $\mu_0\rho\,\partial_x(\delta\mathcal F_\MF/\delta\rho)
= D_0\partial_x\rho + \mu_0\rho\,\partial_x(U*\rho)$, and its noise satisfies
detailed balance with that mobility. Its stationary measure is therefore,
formally, $\propto e^{-\beta\mathcal F_\MF[\rho]}$. The nonlinear terms and the
multiplicative noise are exactly what is needed for this; they do not add
physics beyond $\mathcal F_\MF$.

\paragraph{Equation of state and diffusivity.} For uniform $\rho$, the free
energy density is $f_\MF = k_BT(\rho\ln\rho-\rho) + \tfrac12\Uh(0)\rho^2$, with
$\Uh(0)=\int U\,dx = \sqrt\pi\,\varepsilon\ell$. Hence
\begin{equation}\label{eq:DMF}
  P_\MF = k_BT\rho + \tfrac12\Uh(0)\rho^2, \qquad
  D_\MF = \mu_0\partial_\rho P_\MF = 1 + \beta\Uh(0)\,\rho
        = 1 + \sqrt\pi\,\beta\varepsilon\,\rho\ell ,
\end{equation}
which is the linearized-DK result S103. The mean-field equation of state has
$B_2^\MF = \tfrac12\beta\Uh(0)$ and $B_3^\MF=B_4^\MF=\dots=0$.

\paragraph{Structure factor.} The Gaussian fluctuations of \eqref{eq:FMF} give
the random-phase approximation
\begin{equation}\label{eq:SMF}
  S_\MF(k) = \frac{1}{1+\beta\bar\rho\,\Uh(k)}, \qquad
  S_\MF(0) = \frac{1}{D_\MF}.
\end{equation}
The mean-field DKE is thus internally consistent: it obeys the same
fluctuation--dissipation balance \eqref{eq:S}, but with the mean-field
pressure in place of the true one.

\section{Why the two differ}\label{sec:why}

\paragraph{Mayer function versus potential.} Comparing \eqref{eq:Dvir} with
\eqref{eq:DMF}, mean field replaces the Mayer function by its first-order
expansion,
\begin{equation}
  f_M = e^{-\beta V}-1 \;\longrightarrow\; -\beta V,
\end{equation}
which gives $B_2^\MF = \tfrac12\beta\int V$ in place of
$B_2 = -\tfrac12\int f_M$, and it drops all higher virial coefficients. Since
$e^{-u}-1 > -u$ for $u>0$, a repulsive potential always has $B_2 < B_2^\MF$
(Table~\ref{tab:B}): mean field overestimates the pressure, and with it $D$,
and underestimates $S$.

\paragraph{Direct correlation function.} Equivalently, write the exact
Ornstein--Zernike form $S(k) = 1/(1-\bar\rho\,\hat c(k))$. The RPA
\eqref{eq:SMF} is the closure $c(r) = -\beta V(r)$. At low density the exact
direct correlation function is instead $c(r)\simeq f_M(r)$, which gives
$S(0)\simeq 1/(1+2B_2\bar\rho)$, consistent with \eqref{eq:S} and
\eqref{eq:Dvir} to first order in $\bar\rho$.

\paragraph{Physical picture.} Mean field lets every particle feel the average
density at every separation, i.e.\ it sets the pair correlation $g(r)=1$.
In reality particles avoid each other inside the core, $g(r)\approx
e^{-\beta V(r)}<1$ for $r\lesssim\ell$, so each feels less repulsion than the
average implies. The true pressure rises more slowly with $\rho$, the system
is more compressible, $S$ is larger and collective diffusion is slower.

\paragraph{When mean field is accurate.} The discrepancy is controlled by how
far $e^{-\beta V}-1$ departs from $-\beta V$, which grows with $\beta\varepsilon$.
It also shrinks at high density, where each particle overlaps many
neighbours within $\ell$ and its own pair correlations matter less; this is
the regime $\rho\ell\gtrsim 0.6$ where S103 is found to be accurate at
$\beta\varepsilon=2$. Table~\ref{tab:gap} gives the size of the effect. The
virial series is truncated after $B_4$; the column ``$\rho^3$ term'' shows its
last retained contribution to $D$, a measure of how far it can be trusted.

\begin{table}[h]
\centering
\begin{tabular}{cccccccc}
\toprule
$\beta\varepsilon$ & $\rho\ell$ & $D_\MF$ & $D_\vir$ (to $\rho^2$) & $\rho^3$ term
  & $S_\MF$ & $S_\vir$ & gap in $S$ \\
\midrule
1 & 0.1 & 1.177 & 1.135 & $-0.000$ & 0.849 & 0.881 & 3.6\% \\
1 & 0.4 & 1.709 & 1.615 & $-0.010$ & 0.585 & 0.619 & 5.5\% \\
2 & 0.2 & 1.709 & 1.479 & $+0.008$ & 0.585 & 0.676 & 13\% \\
3 & 0.2 & 2.063 & 1.616 & $+0.027$ & 0.485 & 0.619 & 22\% \\
3 & 0.3 & 2.595 & 2.029 & $+0.091$ & 0.385 & 0.493 & 22\% \\
5 & 0.1 & 1.886 & 1.337 & $+0.008$ & 0.530 & 0.748 & 29\% \\
5 & 0.2 & 2.772 & 1.782 & $+0.065$ & 0.361 & 0.561 & 36\% \\
\bottomrule
\end{tabular}
\caption{Mean-field and virial diffusivity and $S(k\to0)=1/D$ for the 1D
Gaussian core. $S_\vir$ is computed from $D_\vir$ to $\rho^2$.}
\label{tab:gap}
\end{table}

\noindent At $\beta\varepsilon=1$ the gap stays at 4--6\% for all
$\rho\ell\le0.4$: raising the density does not magnify it. At
$\beta\varepsilon=3$, $\rho\ell=0.2$ it is 22\% in $S$ (25\% in $D$) while
the $\rho^3$ term is only 1.7\% of $D$, so including $B_4$ keeps the virial
$D$ accurate. This is the parameter set of the \code{e3r02} inputs. Note also
that the size of the gap is relative to the interaction correction $D-1$, not
to $D$: at $\beta\varepsilon=1$, $\rho\ell=0.1$, mean field overestimates
$D-1$ by about 30\%, but $D-1$ is only 13--18\% of $D$.

\section{Nonlinearity, multiplicative noise and coarse-graining}\label{sec:cg}

The DKE solved on a grid is nonlinear with multiplicative noise, which raises
the question of whether it captures more than its linearization S103. At the
resolutions where it is well defined, it does not.

\paragraph{Gaussian level: independent of the grid.} The mesh kernel is the
cell-pair average of $U$, $U_m = \Delta x^{-2}\int_{\text{cell }0}\int_{\text{cell }m}U(x-y)$.
It satisfies $\sum_m U_m\,\Delta x = \int U\,dx$ exactly, so
$\Uh_{\Delta x}(0)=\Uh(0)$ and the linearized discrete DKE gives
$S(k\to0) = 1/D_\MF$ at every $\Delta x$. Any dependence on $\Delta x$ must come
from the non-Gaussian terms.

\paragraph{Cells much larger than $\ell$.} Two particles in the same cell
interact through $\beta U_0 \approx \beta\varepsilon\sqrt\pi\,\ell/\Delta x$,
e.g.\ $0.009$ at $\beta\varepsilon=1$, $\Delta x=200\ell$, and $0.05$ at
$\beta\varepsilon=3$, $\Delta x=100\ell$. At such weak coupling
$e^{-\beta U_0}-1 = -\beta U_0\,(1+O(\beta U_0))$: the coarse model's own
second virial coefficient is the mean-field one, to within
$\sim\beta U_0/2\lesssim 3\%$ of the interaction correction. The
short-range correlations that distinguish $B_2$ from $B_2^\MF$ live at
separations $\lesssim\ell$, below the grid scale, and are averaged away. The
remaining nonlinear effects are fluctuation corrections of relative size
$\langle\delta\rho^2\rangle/\bar\rho^2 \approx 1/N_{\rm cell}$, with
$N_{\rm cell}=\bar\rho\,\Delta x$ particles per cell. They multiply the
interaction correction $D-1$ and are therefore of order $1\%$ of $D$ for
$N_{\rm cell}=20$. (The ideal-gas part is unaffected by the nonlinearity: with
constant $D_0$ and $\sigma=2\rho$, the equation for
$\langle\delta\rho\,\delta\rho\rangle$ closes exactly and gives $S=1$.)

\paragraph{Cells comparable to or smaller than $\ell$.} When the grid resolves
the core, the nonlinear terms do matter, but they do not reproduce the
particles' exclusion statistics. The field $\rho$ is continuous; its entropy
$\rho\ln\rho$ is Stirling's approximation to $\ln n!$ and is valid only for
$N_{\rm cell}\gg1$. The interaction $\tfrac12\iint\rho U\rho$ includes each
particle's interaction with itself. With $N_{\rm cell}\ll1$, relative
fluctuations are of order $N_{\rm cell}^{-1/2}>1$, so $\rho$ is often near or
below zero, where the noise is switched off. None of this is the Mayer factor
$e^{-\beta V}$.

The runs at $\Delta x=\ell$ ($\beta\varepsilon=1$, $\rho\ell=0.1$, so
$N_{\rm cell}=0.1$) confirm this (Table~\ref{tab:dx1}). The particles follow
the virial prediction, $1/D_\vir=0.881$, but the grid DKE falls below even the
mean-field value: refining the grid moves it further from the exact result.

\begin{table}[h]
\centering
\begin{tabular}{ccccc}
\toprule
$k\ell$ & DKE, $\Delta x=\ell$ & RPA $1/(1+\beta\bar\rho\Uh(k))$ & particles & $1/D_\vir$ ($k\to0$) \\
\midrule
0.105 & 0.799 & 0.850 & 0.877 & 0.881 \\
0.628 & 0.818 & 0.862 & 0.893 & \\
1.15  & 0.857 & 0.889 & 0.924 & \\
2.72  & 0.970 & 0.973 & 0.995 & \\
\bottomrule
\end{tabular}
\caption{$F(k,0)$ measured at $\beta\varepsilon=1$, $\rho\ell=0.1$ on a grid
with $\Delta x=\ell$ (0.1 particles per cell), for the grid DKE and the
particles, against the mean-field (RPA) value with
$\Uh(k)=\sqrt\pi\,\varepsilon\ell\,e^{-k^2\ell^2/4}$.}
\label{tab:dx1}
\end{table}

\noindent The coarse-grained DKE is therefore well defined only with many
particles per cell, and in that regime it is a mean-field theory at the scale
of the interaction: its long-wavelength $S$ and $D$ are $1/D_\MF$ and
$D_\MF$ for any $\Delta x$. Beyond-mean-field thermodynamics must be supplied
explicitly, either through the exact $D(\rho)=\mu_0\partial_\rho P$ in MFT, or
by the particle dynamics itself. The exact Dean equation for the empirical
density $\sum_i\delta(x-x_i)$ does contain the correct physics, because it is
exact; it is the coarse-graining to a smooth field on cells larger than $\ell$,
with a smoothed kernel, that reduces it to mean field.

\section{Consequences for the ISF comparison}

For the \code{e3r02} runs ($\beta\varepsilon=3$, $\rho\ell=0.2$, $\Delta x=100\ell$,
20 particles per cell), \eqref{eq:S} and \eqref{eq:F} give the predictions of
Table~\ref{tab:pred}, with $\keff^2=(4/\Delta x^2)\sin^2(k\Delta x/2)$ in the
decay rate of the grid models.

\begin{table}[h]
\centering
\begin{tabular}{lccc}
\toprule
model & $D$ & $S(k\to0)=F(k,0)$ & decay rate \\
\midrule
particles (exact)            & $D_\vir = 1.643$ & 0.609 & $D_\vir k^2$ \\
MFT, virial $D(\rho)$        & 1.643 & 0.609 & $D_\vir\keff^2$ \\
MFT, mean-field $D(\rho)$    & 2.063 & 0.485 & $D_\MF\keff^2$ \\
coarse-grained DKE           & $D_\MF = 2.063$ & 0.485 & $D_\MF\keff^2$ \\
\bottomrule
\end{tabular}
\caption{Predicted equilibrium statics and dynamics at $\beta\varepsilon=3$,
$\rho\ell=0.2$.}
\label{tab:pred}
\end{table}

\noindent Agreement between the DKE and MFT with the mean-field $D$ checks the
consistency of the two solvers; it is not a test of physics, since both share
the same linear theory. The physics test is whether the particles follow the
virial row (MFT) or the mean-field row (DKE).

\section{Correcting the DKE: an effective interaction kernel}\label{sec:ceff}

\paragraph{Not a renormalized $D_0$.} The kinetic coefficients of the DKE are
already exact: its mobility $\mu_0\rho$ is the exact $\sigma/(2k_BT)$
\eqref{eq:sigma}, and coarse-graining does not change it. Rescaling $D_0$
consistently (Fickian term, drift mobility and noise) only changes the time
scale: $D_c = D_0(1+\beta\bar\rho\Uh(0))$ and $\sigma = 2D_0\bar\rho$ give
$S = 1/(1+\beta\bar\rho\Uh(0))$ for any $D_0$. Changing only the Fickian term,
to $D_0' = D_\vir - \beta\bar\rho\Uh(0)$, reproduces $D_\vir$ at the linear
level, but $D_0'$ must depend on $\rho$, it breaks detailed balance with
$\mathcal F_\MF$, and it amounts to the MFT with the true $D(\rho)$. The error
is in the free energy, so the correction belongs in the interaction.

\paragraph{Effective kernel.} Replace $V$ in \eqref{eq:dk} by
\begin{equation}\label{eq:ueff}
  U_{\rm eff}(r) = -k_BT\,c(r),
\end{equation}
with $c$ the direct correlation function, keeping the mobility $D_0\rho$. The
Gaussian fluctuations then give the exact Ornstein--Zernike structure factor
and the exact collective diffusivity,
\begin{equation}
  S(k) = \frac{1}{1-\bar\rho\,\hat c(k)}, \qquad
  D_c = D_0\big(1-\bar\rho\,\hat c(0)\big) = \frac{D_0}{S(0)} = \mu_0\partial_\rho P ,
\end{equation}
the standard dynamical-DFT (Ramaswamy--Yussouff) form. The low-density
closure $c\simeq e^{-\beta V}-1$ gives $\beta\hat U_{\rm eff}(0) = 2B_2$, i.e.
$D = 1+2B_2\rho$ (1.475 at $\beta\varepsilon=3$, $\rho\ell=0.2$), still 10\%
below $D_\vir$ because $B_3$ and $B_4$ matter there. Since the grid only
resolves $k\ell\ll1$ ($\Delta x = 100\ell$), where $\hat c(k)\simeq\hat c(0)$,
it suffices to match the compressibility: keep the Gaussian shape and choose
its strength so that $1+\sqrt\pi\,\beta\varepsilon_{\rm eff}\,\rho\ell = D_\vir$,
\begin{equation}
  \beta\varepsilon_{\rm eff} = \frac{D_\vir-1}{\sqrt\pi\,\rho\ell}
   = \frac{0.6428}{\sqrt\pi\times0.2} = 1.813
   \qquad(\text{bare } \beta\varepsilon = 3).
\end{equation}
This is \code{DK/inputs\_fv\_gauss\_1d\_ensemble\_e3r02\_ceff}
(\code{amr.ip\_eps} $=1.813\,N$), otherwise identical to the bare-kernel DK run.

\paragraph{Results.} Table~\ref{tab:ceff} collects the measured collective
diffusivity, fitted from the decay of $F(k,\tau)$, and $F(k,0)$, averaged over
the modes $q=8,12,\dots,28$ ($k\ell\le0.03$). For the grid runs the fit is the
slope of $\ln F$ divided by $\keff^2$. For the particles, whose mesh density is
a nearest-grid-point particle count, the $q$ mode contains the aliases
$k_m = k+2\pi m/\Delta x$ with weights $\mathrm{sinc}^2(k_m\Delta x/2)$, each
decaying at $Dk_m^2$, and $D$ is fitted to
$F_p = A\sum_m \mathrm{sinc}^2(k_m\Delta x/2)\,e^{-Dk_m^2\tau}$.

\begin{table}[h]
\centering
\begin{tabular}{lcccc}
\toprule
run & $D$ & $S=F(k,0)$ & $D\,S$ & theory $D$, $S$ \\
\midrule
DKE, bare kernel ($\beta\varepsilon=3$)        & $2.065\pm0.006$ & $0.484\pm0.001$ & 0.999 & 2.063, 0.485 \\
MFT, mean-field $D(\rho)$                      & $2.071\pm0.005$ & $0.484\pm0.001$ & 1.002 & 2.063, 0.485 \\
DKE, effective kernel ($\beta\varepsilon_{\rm eff}=1.813$) & $1.642\pm0.007$ & $0.606\pm0.002$ & 0.994 & 1.643, 0.609 \\
MFT, virial $D(\rho)$                          & $1.645\pm0.006$ & $0.607\pm0.001$ & 0.998 & 1.643, 0.609 \\
particles                                      & $1.63\pm0.05^{\,a}$ & $0.606\pm0.011$ & 0.95$^{\,b}$ & 1.643, 0.609 \\
\bottomrule
\end{tabular}
\caption{Measured collective diffusivity and $S(k\to0)$ at
$\beta\varepsilon=3$, $\rho\ell=0.2$, 20 particles per cell. Grid runs: mean
and standard deviation over the six modes. Particles: mean and standard error.
$^a$ modes $q\ge16$, which the particle run follows over several relaxation
times; over all modes $1.57\pm0.05$, the low-$q$ fits being limited by a run of
only 1--3 relaxation times with 77 time origins. $^b$ over all modes.}
\label{tab:ceff}
\end{table}

\noindent Four conclusions follow.
\begin{enumerate}[nosep]
  \item The effective-kernel DKE agrees with the virial MFT within 0.2\% in
        both $D$ and $S$, at every mode. The mean-field error of the DKE is
        therefore entirely thermodynamic, and is removed by \eqref{eq:ueff};
        $D_0=1$ needs no change.
  \item The nonlinear and multiplicative-noise terms of the DKE add nothing
        measurable beyond its linear theory: $D\,S-1$ stays below 1\% for all
        grid runs, consistent with the $O(1/N_{\rm cell})$ estimate of
        Sec.~\ref{sec:cg}.
  \item The particles follow the virial (and effective-kernel) prediction and
        exclude mean field: $S$ is $11$ standard errors from $1/D_\MF$ and $D$
        (modes $q\ge16$) $8$ standard errors from $D_\MF$.
  \item Matching only $\hat c(0)$ suffices for all resolved modes, because
        $\hat c(k)$ is flat for $k\ell\ll1$.
\end{enumerate}

\paragraph{Limitations.} $\beta\varepsilon_{\rm eff}$ was fixed from the
virial equation of state at this density: it is a one-parameter matching of
the compressibility, not a prediction. Since $c$ depends on $\rho$, the kernel
must be recomputed at another density, and a single kernel is not exact for an
inhomogeneous $\rho(x)$; the MFT, whose $D(\rho)$ depends on the local
density, handles both automatically. Resolving structure at $k\ell\sim1$
would require the full $\hat c(k)$, not only $\hat c(0)$. Finally, the
particle statistics ($\pm2$--$3\%$) cannot resolve corrections at the 1\%
level, such as those of the truncated virial series (the $B_4$ term is 1.7\%
of $D$).

\section{Summary}

\begin{enumerate}[nosep]
  \item For Brownian particles $\sigma=2\mu_0k_BT\rho$ exactly, so
        $S(k\to0)=\sigma/(2\rho D)=1/D$, and $D=\mu_0\partial_\rho P$ is the
        inverse compressibility: statics and collective dynamics are both set by
        the equation of state.
  \item MFT with $D$ from the virial expansion uses the exact equation of
        state, including pair correlations through the Mayer function
        $e^{-\beta V}-1$.
  \item The coarse-grained DKE is the gradient flow of the mean-field free
        energy \eqref{eq:FMF}. Its equation of state has
        $B_2^\MF=\tfrac12\beta\Uh(0)$ and no higher coefficients, i.e.\ it sets
        $g(r)=1$. It overestimates the pressure and $D$ and underestimates $S$.
  \item Its nonlinearity and multiplicative noise only make its stationary
        measure $e^{-\beta\mathcal F_\MF}$; with many particles per cell their
        corrections are $O(1/N_{\rm cell})$ of the interaction correction.
        Refining the grid does not recover the virial result, and at
        $N_{\rm cell}\ll1$ the field theory itself is not reliable.
  \item The gap grows with $\beta\varepsilon$, and is about 22\% in $S$ at
        $\beta\varepsilon=3$, $\rho\ell=0.2$, where the virial $D$ is still
        accurate.
  \item The fix is in the interaction, not in $D_0$: with the effective
        kernel $-k_BT\,c(r)$, matched to the compressibility
        ($\beta\varepsilon_{\rm eff}=1.813$ for $\beta\varepsilon=3$), the DKE
        reproduces the virial MFT within 0.2\%, and both agree with the
        particle simulations within their statistics (Table~\ref{tab:ceff}).
\end{enumerate}

\end{document}
